\documentclass[12pt, a4paper]{article}

\usepackage{amsmath, amssymb}
\usepackage{fontspec}
\setmainfont{FreeSerif}
\usepackage{geometry}
\geometry{margin=1in}
\usepackage[colorlinks=true, linkcolor=blue, urlcolor=blue, citecolor=blue]{hyperref}

\setcounter{secnumdepth}{0} 

\begin{document}

\begin{center}
    \Huge \textbf{Անկյունագծայնացում և SVD}
\end{center}

\tableofcontents
\vspace{2em}
\hrule
\vspace{2em}

\section{1. Մատրիցների անկյունագծայնացում}

\subsection{1.1 Սահմանում}

Թող \(A\) լինի քառակուսի մատրից։ \(A\) մատրիցը կոչվում է \textbf{անկյունագծայնացվող}, եթե գոյություն ունի հակադարձելի \(P\) մատրից և անկյունագծային \(D\) մատրից այնպես, որ

$$
A = P D P^{-1}.
$$

Այստեղ \(D\)-ն անկյունագծային մատրից է, իսկ \(P\)-ն հակադարձելի մատրից։

\vspace{1em}
\hrule
\vspace{1em}

\subsection{1.2 Մատրիցի աստիճանի հաշվում}

Անկյունագծայնացումը թույլ է տալիս էականորեն պարզեցնել մատրիցի բարձր աստիճանների հաշվարկը։

Օրինակ՝

$$
A^2 = (PDP^{-1})(PDP^{-1}) 
= P D (P^{-1}P) D P^{-1} 
= P D^2 P^{-1}.
$$

Ընդհանուր դեպքում ստացվում է

$$
A^n = P D^n P^{-1}.
$$

Քանի որ \(D\)-ն անկյունագծային մատրից է, ապա \(D^n\)-ը հաշվվում է շատ հեշտ՝ յուրաքանչյուր անկյունագծային տարրը բարձրացնելով \(n\)-րդ աստիճանի։

\vspace{1em}
\hrule
\vspace{1em}

\subsection{1.3 Սիմետրիկ մատրիցներ}

Եթե \(A\) մատրիցը սիմետրիկ է, այսինքն

$$
A = A^T,
$$

ապա այն անկյունագծայնացվում է \textbf{օրթոգոնալ մատրիցի} միջոցով։

Այսինքն գոյություն ունի օրթոգոնալ \(P\) մատրից այնպես, որ

$$
A = P D P^T.
$$

Քանի որ օրթոգոնալ մատրիցի համար

$$
P^{-1} = P^T,
$$

ստացվում է ավելի պարզ ներկայացում։

\vspace{1em}
\hrule
\vspace{1em}

\subsection{1.4 Հակադարձ մատրիցի սեփական արժեքները}

Թող

$$
Ax = \lambda x
$$

լինի \(A\) մատրիցի սեփական արժեքի և սեփական վեկտորի հավասարումը, որտեղ \(\lambda \neq 0\)։

Այդ դեպքում \(A^{-1}\) մատրիցի սեփական արժեքը կլինի

$$
\frac{1}{\lambda}.
$$

\subsubsection{Ապացույց}

$$
Ax = \lambda x
$$

կիրառենք \(A^{-1}\) երկու կողմերին․

$$
A^{-1}(Ax) = A^{-1}(\lambda x)
$$

ստացվում է

$$
Ix = \lambda A^{-1}x
$$

այսինքն

$$
A^{-1}x = \frac{1}{\lambda}x.
$$

\vspace{1em}
\hrule
\vspace{1em}

\section{2. Մատրիցի ռանգ (Rank)}

\subsection{2.1 Սահմանում}

Թող \(A \in M_{m\times n}\) լինի մատրից։

Մատրիցի \textbf{տողային ռանգը} սահմանվում է որպես այդ մատրիցի գծորեն անկախ տողերի առավելագույն քանակը։

\vspace{1em}
\hrule
\vspace{1em}

\subsection{2.2 Հիմնական թեորեմ}

Կամայական մատրիցի դեպքում տողային և սյունային ռանգերը համընկնում են։

$$
\text{rank}_{row}(A) = \text{rank}_{col}(A) = \text{rank}(A).
$$

Բացի այդ, միշտ ճիշտ է նաև

$$
\text{rank}(A) = \text{rank}(A^T).
$$

\vspace{1em}
\hrule
\vspace{1em}

\section{3. Եզակի արժեքների վերլուծություն (SVD)}

\subsection{3.1 Թեորեմ}

Թող \(A \in M_{m\times n}\) լինի ցանկացած մատրից։ Այդ դեպքում գոյություն ունեն

\begin{itemize}
    \item \(U \in M_{m\times m}\) օրթոգոնալ մատրից,
    \item \(V \in M_{n\times n}\) օրթոգոնալ մատրից,
    \item \(\Sigma \in M_{m\times n}\) անկյունագծային տեսքի մատրից,
\end{itemize}

այնպես, որ

$$
A = U \Sigma V^T.
$$

Այստեղ

\begin{itemize}
    \item \(U\)-ի սյուները \(AA^T\) մատրիցի սեփական վեկտորներն են,
    \item \(V\)-ի սյուները \(A^TA\) մատրիցի սեփական վեկտորներն են,
    \item \(\Sigma\)-ի անկյունագծային տարրերը կոչվում են \textbf{եզակի արժեքներ}։
\end{itemize}

Եթե \(\lambda_i\)-ները \(A^TA\)-ի սեփական արժեքներն են, ապա

$$
\sigma_i = \sqrt{\lambda_i}.
$$

\vspace{1em}
\hrule
\vspace{1em}

\subsection{3.2 Երկրաչափական մեկնաբանություն}

SVD-ն ներկայացնում է գծային ձևափոխությունը երեք հաջորդական գործողությունների միջոցով․

\begin{enumerate}
    \item պտույտ \(V^T\)
    \item ձգում \(\Sigma\)
    \item կրկին պտույտ \(U\)
\end{enumerate}

Միավոր շրջանագիծը գծային ձևափոխության արդյունքում վերածվում է էլիպսի, որի առանցքները տրված են \(\sigma_i u_i\) վեկտորներով։

\vspace{1em}
\hrule
\vspace{1em}

\section{4. Գործնական օրինակների լուծում}

\subsection{Խնդիր 256. \(A^{10}\)-ի հաշվում}

Տրված է

$$
A =
\begin{pmatrix}
1 & 0 \\
-1 & 2
\end{pmatrix}.
$$

\textbf{Սեփական արժեքներ}

$$
\det(A-\lambda I) = (1-\lambda)(2-\lambda)=0
$$

$$
\lambda_1 = 1, \qquad \lambda_2 = 2
$$

\textbf{Անկյունագծայնացում}

$$
A =
\begin{pmatrix}
1 & 0 \\
1 & 1
\end{pmatrix}
\begin{pmatrix}
1 & 0 \\
0 & 2
\end{pmatrix}
\begin{pmatrix}
1 & 0 \\
-1 & 1
\end{pmatrix}
$$

\textbf{Աստիճանի հաշվարկ}

$$
A^{10} =
\begin{pmatrix}
1 & 0 \\
-1023 & 1024
\end{pmatrix}
$$

\vspace{1em}
\hrule
\vspace{1em}

\section{5. Կիրառություն. Կեղծ հակադարձ մատրից}

Եթե

$$
Ax=b
$$

համակարգը լուծում չունի, ապա փնտրում ենք այնպիսի \(x_0\), որի համար

$$
\|Ax_0 - b\|
$$

լինի նվազագույն։

Լուծումը

$$
x_0 = A^+ b
$$

\vspace{1em}
\hrule
\vspace{1em}

\section{Մուր–Պենրոուզի կեղծ հակադարձ մատրից}

Եթե մատրիցը քառակուսի չէ կամ հակադարձելի չէ, օգտագործվում է \textbf{կեղծ հակադարձ մատրիցը}։

Եթե

$$
A \in M_{m\times n}
$$

ապա

$$
A^+ \in M_{n\times m}
$$

Եթե

$$
A = U\Sigma V^T
$$

ապա

$$
A^+ = V \Sigma^+ U^T
$$

որտեղ \(\Sigma^+\)-ը ստացվում է \(\Sigma\)-ից՝ անկյունագծային տարրերը փոխարինելով իրենց հակադարձներով։

\vspace{1em}
\hrule
\vspace{1em}

\section{Կիրառություն. Լավագույն մոտարկման խնդիրը}

Եթե \(Ax=b\) համակարգը ճշգրիտ լուծում չունի, ապա փնտրում ենք

$$
\|Ax_0 - b\| \to \min
$$

Լավագույն լուծումը

$$
x_0 = A^+ b
$$

\vspace{1em}
\hrule
\vspace{1em}

\section{Գործնական օրինակի լրիվ լուծում}

Տրված է

$$
A =
\begin{pmatrix}
1 & 1 \\
0 & 1 \\
1 & 0
\end{pmatrix},
\quad
b =
\begin{pmatrix}
1 \\
2 \\
3
\end{pmatrix}
$$

\textbf{Լուծում}

$$
x_0 =
\begin{pmatrix}
5/3 \\
2/3
\end{pmatrix}
$$

Սխալը

$$
\|Ax_0 - b\|^2 = \frac{16}{3}
$$

որը նվազագույն հնարավոր սխալն է։

\end{document}